\section{Implementación de la comunicación mediante gRPC}
\label{sec:grpc}

Cada nodo mantiene un servidor gRPC y un cliente dentro del mismo proceso.
El servidor atiende RPC mediante un grupo de cuatro hilos y escucha en
\texttt{127.0.0.1}. El cliente utiliza un canal por destino y puerto.
\texttt{nodo.py} integra ambos componentes con el reloj de Lamport y el
registro CSV de \texttt{src/comun/}, sin depender del transporte TCP.

\subsection{Contrato y generación}

El servicio \texttt{Mensajeria} expone la operación unaria
\texttt{EnviarMensaje}: cada solicitud produce un acuse. El contrato
de \texttt{proto/mensajes.proto} contiene los mismos datos lógicos que TCP:

\begin{verbatim}
syntax = "proto3";
package pe_u1;
service Mensajeria {
  rpc EnviarMensaje(Mensaje) returns (Acuse);
}
message Mensaje {
  string sender = 1;
  int64 timestamp = 2;
  string message = 3;
  string message_id = 4;
}
message Acuse {
  string sender = 1;
  int64 timestamp = 2;
  string message = 3;
  string message_id = 4;
  string ack_for = 5;
  bool ok = 6;
}
\end{verbatim}

Los campos de texto usan \texttt{string} y el reloj usa un entero de
64 bits con signo. El servicio rechaza relojes negativos e identificadores
o emisores vacíos mediante el estado \texttt{INVALID\_ARGUMENT}. Los números
de campo identifican los valores en la representación binaria. El acuse
confirma el identificador original mediante \texttt{ack\_for}; su propio
identificador añade el sufijo \texttt{-ack}.

La generación sigue el procedimiento documentado por
gRPC\footnote{\texttt{https://grpc.io/docs/languages/python/basics/}}.
Desde la raíz, en \texttt{cmd.exe}, se ejecuta:

\begin{verbatim}
python -m grpc_tools.protoc -Isrc/grpc/generado=proto ^
  --python_out=. --grpc_python_out=. proto/mensajes.proto
\end{verbatim}

El prefijo virtual determina la ruta \texttt{src/grpc/generado/} y las
importaciones del paquete. Allí se generan \texttt{mensajes\_pb2.py}
y \texttt{mensajes\_pb2\_grpc.py}, sin modificaciones manuales.

Se comprobaron las siguientes versiones: Python 3.12, grpcio 1.85.3,
grpcio-tools 1.85.3 y protobuf 7.36.2.

\subsection{Reloj, acuses y latencia}

Cada proceso comparte un único reloj entre cliente y servidor.
Antes de enviar la solicitud se incrementa el contador. El servidor
aplica \(\max(local, recibido)+1\) al recibirla y vuelve a incrementarlo
al enviar el acuse. El cliente valida el emisor, los identificadores,
el indicador de éxito y que el reloj del acuse sea posterior al de su
solicitud; después registra la recepción. Un bloqueo agrupa cada avance
del reloj con su escritura CSV para conservar el orden local entre hilos.

La conexión inicial se establece antes del intervalo de latencia y el
canal se reutiliza en las solicitudes posteriores. La medición incluye
la serialización de la solicitud, el RPC, la deserialización del acuse
y su validación, y termina antes del registro de la recepción:

\begin{verbatim}
inicio = perf_counter()
acuse = stub.EnviarMensaje(solicitud, timeout=restante,
                          wait_for_ready=False)
\end{verbatim}

Después de validar el acuse se calcula
\texttt{(perf\_counter() - inicio) * 1000}. El tiempo máximo de la
operación completa, incluida la conexión inicial, es de cinco segundos
por defecto. El RPC recibe el tiempo restante de ese presupuesto.
\texttt{--timeout} permite valores mayores que cero y de hasta diez
segundos, inferiores a la espera de quince segundos del ejecutor.
Tras un fallo se cierra el canal afectado. La aplicación no reintenta
automáticamente la solicitud, porque el servidor podría haberla procesado
antes de agotarse el plazo. Una orden posterior puede abrir otro canal.

\subsection{Interfaz de ejecución y errores}

Los argumentos, las órdenes JSON por entrada estándar y las respuestas
son compatibles con el ejecutor existente. Una respuesta inicial con
\texttt{ok=true} anuncia que el servidor está activo. Las órdenes contienen
\texttt{destino}, \texttt{puerto}, \texttt{message} y \texttt{message\_id}.
Una respuesta exitosa devuelve el identificador, \texttt{ok=true} y
\texttt{latencia\_ms}. Los fallos devuelven \texttt{ok=false} y una
descripción; los errores gRPC incluyen el estado y el detalle del RPC.
La salida estándar contiene exclusivamente respuestas JSON. Los
diagnósticos se reservan para la salida de error. Al terminar la entrada
estándar, el proceso cierra el cliente y el servidor.

Para iniciar A desde la raíz se utiliza:

\begin{verbatim}
python -m src.grpc.nodo --nombre nodo-A --puerto 5000 ^
  --corrida prueba-byron --eventos data/grpc/prueba-A.csv
\end{verbatim}

B y C se inician en terminales distintas con sus nombres y puertos
5001 y 5002, y con rutas CSV propias. Los archivos de eventos deben ser
nuevos. Cuando se omite \texttt{--eventos}, el reloj sigue activo pero
no se escriben eventos, conforme al protocolo de medición.

\subsection{Comprobación manual en el equipo de Byron}

El 21 de noviembre de 2026, Byron inició los tres nodos en terminales
de su equipo y ejecutó tres solicitudes: \texttt{B-001} de A hacia B,
\texttt{B-002} de B hacia C y \texttt{B-003} de C hacia A. Las tres
recibieron un acuse válido. Después cerró los procesos mediante EOF.
Las capturas originales de los intercambios y del cierre se conservan
en \texttt{evidencias/byron/}, con usuario, equipo, fecha y zona horaria
visibles.

Los tres CSV de \texttt{data/grpc/practica-byron-01/} contienen doce
eventos: envío y recepción de tres solicitudes y sus tres acuses.
El primer intercambio registra los valores 1, 2, 3 y 4, correspondientes
al envío de A, recepción en B, envío del acuse de B y recepción en A.
En los seis pares registrados, el reloj de recepción es mayor que el
de envío. Los relojes locales también crecen estrictamente.
Los eventos abarcan desde las 18:49:51 hasta las 18:52:27, UTC-05:00.
Después del cierre se calcularon y verificaron las sumas de los CSV
en el archivo \texttt{SHA256SUMS} de esa misma carpeta.

Esta comprobación manual documenta el intercambio entre los tres
procesos. Por su número de envíos y sus condiciones, no sustituye
las repeticiones, tamaños y concurrencia de la campaña experimental,
ni permite concluir qué transporte tiene menor latencia.

\subsection{Alcance de las comprobaciones y asistencia de IA}

Codex, agente de OpenAI basado en GPT-6, generó el contrato, los módulos
gRPC, las pruebas automatizadas y
esta sección, reutilizando las utilidades comunes y el ejecutor existentes.
Codex ejecutó comprobaciones técnicas con conexiones locales reales y
datos temporales: acuses, plazos, errores, reutilización del canal,
relojes, concurrencia e integración de los tres procesos. Estas
comprobaciones no constituyen la campaña experimental ni evidencias
personales de ejecución de Byron. La práctica manual descrita arriba
fue operada por Byron siguiendo la orientación de Codex. Codex comprobó
los registros, copió sin editar las capturas proporcionadas por Byron,
verificó sus copias y calculó las sumas de los CSV cerrados. La redacción
de esta comprobación también fue asistida por Codex.
